\section{讨论：边界伪像、退化谱与解释层对接}

\subsection{边界伪像的组合学来源}

分裂 $21=18\oplus 3$ 是线性窗口语义与周期贴合语义不一致的必然结果。该结论不依赖任何动力学选择，仅由禁词约束与“首尾相邻”的语义差异决定。

\subsection{退化谱的结构性意义}

定理 \ref{thm:deg-spectrum-6} 表明：在最小核 $m=6$ 中，纤维退化度由尾部隐藏数字可行性与上界约束共同决定，而非经验现象。该机制提供一种离散原型：可见结构由规范形决定，不可见结构由尾部自由度的可行组合数决定。

\subsection{解释层对接（不进入定理层）}

在高维周期—低维投影的解释框架中，可见坐标与内部坐标分离；本文的纤维 $\mathcal{F}_m(x)$ 可作为内部标签的离散原型，闭包固定点可视为“观测一致性”要求下的现实集合自洽条件。

\subsection{分辨率无关的正确表述（解释层用语纪律）}

本文中参数 $m$ 是 Zeckendorf 窗口长度，改变 $m$ 会改变有限输出空间 $X_m$ 的规模与纤维退化统计，因此不应以“不同 $m$ 得到同一原始输出”表述跨尺度关系。跨尺度“同一性”的严格口径是相容性（coherence）：高分辨率输出经由自然 restriction 投影后复现低分辨率输出（见引理 \ref{lem:foldm-pointwise-compat}），从而低分辨率结构作为高分辨率结构的因子被保留；更大 $m$ 仅提供更长前缀与更细的条件结构。

